\section{Introduction}
\label{sec:introduction}

\subsection{Motivation} 
Aritficial neural networks have shown great results
and improvements in the last years, impacting fields such 
as science, poltics and philosophy. Large language models, image detection,
drones and autonomous systems, all have variations of 
artificial neural networks as its basis. In politics we can hear
about intelligent machines, aritificial intelligence and super intelligence.
The motivation for this project is to learn about neural networks, demistify
its supernatural appearance, and understand the mechanics
and its capacity to learn and solve known functions and data relations.  
The method and results will be compared to known methods of linear
regression and classification. As a test of strength and calibration.
% Dette må enten videreutvikles eller fjernes.


\subsection{What i have done}
\subsubsection{Regression problem}
We have created a feed forward neural network; FFNN, used stochastic gradient descent; SGD,
to find the minimum of the cost function, which is the MSE for the regression problem. 
The gradients are calculated using backpropagation, which uses the kernal rule of derivation
on the FFNN in reverse; from the output and target MSE derivative to the
derivative of the parameters of the first layer. 
The gradients are compared with jax, automatic differentiation,
to ensure agreement. In the SGD the descent method used on the parameters and gradients
is the Plain gradient descent; regular updating, 
RMSprop; complex updating, Adam; best updating.  
Different layer width and depths are explored, along with the activation functions
in the hidden layers, different learning rates are explored and compared. 
The output layer has the identity as activaiton function.
The choice of activation function is the identity since the problem of choice is a 
regression problem. This network is compared with OLS for the one dimensional runge function.
Further L1 and L2 penalties are applied to the cost function, with a variable
penalty value $\lambda$, and these models are compared with Lasso and Ridge.  
\subsubsection{Classification problem}
For the classification problem the difference is the output activation function
and the cost function. The output function is the softmax function, which..
The cost function is the cross entropy, and the derivative of this function
changes the backpropagation algorithm accordingly. This and the data set and following
problem is the difference between this network and the regression network.
The study problem is the multiclass MNIST problem of handwritten digits 
(LeCun, Cortes and Burges), the description of the data set is found at 
https://www.tensorflow.org/datasets/catalog/mnist. 
For the classification network the explored parameters
are whidth of the hidden layer, the learning rate $\gamma$, 
the gradient methods plain, rmsprop and adam. Also
the L1 and L2 penalties on the cost function are applied
with different penalty values $\lambda$.  
%QUESTION: Since the cost function is different, how does
%the penalty affect this? Am i misunderstanding something. 

\subsection{In this report}
\ref{sec:theory}
\ref{sec:methods}
%FFNN
%Regression
%Classification
\ref{sec:results}
\ref{sec:discussion}
\ref{sec:conclusion}


